\documentclass[twocolumn]{aastex631}
% Auto-generated by jansky_research.hiblend.write_paper -- do not edit.
% hbReal* come from results/hiblend_metrics.json, hiblend_diagnostics.json,
% hiblend_v2_metrics.json and hiblend_referee{1,2,3}.json (real data only).
\newcommand{\hbSource}{FASHI DR2 (arXiv:2606.31539) x ALFALFA alpha.100 (Haynes+2018, J/ApJ/861/49)}
\newcommand{\hbRealNFashi}{156{,}269}
\newcommand{\hbRealNAlfalfa}{31{,}500}
\newcommand{\hbRealNTargets}{22{,}889}
\newcommand{\hbRealNIsolated}{14{,}375}
\newcommand{\hbRealNPrimary}{2{,}907}
\newcommand{\hbRealNNull}{4{,}560}
\newcommand{\hbRealBeamFast}{2.9}
\newcommand{\hbRealBeamAlfa}{3.5}
\newcommand{\hbRealSignal}{0.045}
\newcommand{\hbRealPowerFrac}{100}
\newcommand{\hbRealInjOne}{1.25}
\newcommand{\hbRealInjOneSd}{0.04}
\newcommand{\hbRealInjZero}{-0.01}
\newcommand{\hbRealChance}{0.39}
\newcommand{\hbRealBetaNull}{0.26}
\newcommand{\hbRealBetaNullErr}{0.09}
\newcommand{\hbRealBetaNullSig}{3.0}
\newcommand{\hbRealBeta}{0.38}
\newcommand{\hbRealBetaErr}{0.14}
\newcommand{\hbRealDzeroTop}{0.039}
\newcommand{\hbRealDzeroTopSig}{8.9}
\newcommand{\hbRealDzeroMid}{-0.025}
\newcommand{\hbRealMatchedNull}{0.12}
\newcommand{\hbRealMatchedNullErr}{0.08}
\newcommand{\hbRealTercLo}{0.38}
\newcommand{\hbRealTercLoErr}{0.12}
\newcommand{\hbRealTercHi}{-0.03}
\newcommand{\hbRealTercHiErr}{0.20}
\newcommand{\hbRealDensity}{-0.001}
\newcommand{\hbRealDensityErr}{0.009}
\newcommand{\hbRealVoneMax}{0.054}
\newcommand{\hbRealVtwoMaxA}{0.015}
\newcommand{\hbRealVtwoMaxB}{0.027}
\newcommand{\hbRealVtwoPA}{0.006}
\newcommand{\hbRealVtwoPB}{0.001}
\newcommand{\hbRealVtwoLoMin}{-0.027}
\newcommand{\hbRealVtwoLoMax}{-0.015}
\newcommand{\hbRealVtwoHiMin}{0.009}
\newcommand{\hbRealVtwoHiMax}{0.015}
\newcommand{\hbRealVtwoSigMin}{2.5}
\newcommand{\hbRealVtwoSigMax}{3.5}
\newcommand{\hbRealVoneMaxB}{0.052}
\newcommand{\hbRealAmpPass}{0.33}
\newcommand{\hbRealMeanPA}{0.36}
\newcommand{\hbRealMeanPB}{0.043}
\newcommand{\hbRealMisfitNull}{0.003}
\newcommand{\hbRealMisfitNullErr}{0.018}
\newcommand{\hbRealRnullMed}{0.0034}
\newcommand{\hbRealRprimMed}{0.0021}
\newcommand{\hbRealMisfitNullMed}{-0.07}
\newcommand{\hbRealConeProdSd}{0.13}
\newcommand{\hbRealMeanMaxB}{0.011}
\newcommand{\hbRealMeanMaxA}{0.009}
\newcommand{\hbRealScaledFrozen}{1.00}
\newcommand{\hbRealScaledProd}{1.54}
\newcommand{\hbRealScaledSd}{0.10}
\newcommand{\hbRealConeProdMin}{0.72}
\newcommand{\hbRealConeProdMax}{1.54}
\newcommand{\hbRealConeFrozMin}{0.18}
\newcommand{\hbRealConeFrozMax}{1.00}
\newcommand{\hbRealVoneMaxMean}{0.051}
\newcommand{\hbRealConeProd}{1.24}
\newcommand{\hbRealConeFrozen}{0.70}
\newcommand{\hbRealConeFrozenSd}{0.13}
\newcommand{\hbRealStripSigMin}{2.6}
\newcommand{\hbRealStripSigMax}{2.9}

\newcommand{\code}[1]{\texttt{#1}}
\shorttitle{An inconclusive two-beam blending test}
\shortauthors{Barbere}
\begin{document}
\title{An Inconclusive Two-Beam FASHI--ALFALFA Test for H\,I Beam Blending}
\author[0009-0008-3289-4447]{Joseph Barbere}
\affiliation{Independent researcher}

\begin{abstract}
Beam blending should raise a galaxy's H\,I flux more in the Arecibo beam
(\hbRealBeamAlfa~arcmin) than in the FAST beam (\hbRealBeamFast~arcmin). We tested this on
\hbRealNTargets\ galaxies matched between FASHI DR2 and ALFALFA $\alpha$.100, with controls
fixed in advance. Power, injection and chance-match controls passed; neighbours too far apart
in velocity to blend gave a slope of $\hbRealBetaNull\pm\hbRealBetaNullErr$ instead of
zero, so the test is inconclusive. The calibration's signal-to-noise misfit predicts only
$\hbRealMisfitNull\pm\hbRealMisfitNullErr$ of that slope, and the failure survives other bootstrap
units; its cause is unknown. A second attempt stopped at its held-out calibration
gate, but that gate judged medians against a tolerance that a perfect calibration would fail
two times in three. In the mean, its calibration passes. No blending slope is quoted.
\end{abstract}

\keywords{H I line emission (690) --- Sky surveys (1464) --- Flux calibration (544) ---
Galaxy environments (2029)}

\section{Test}
\label{sec:test}
In FASHI DR2 \citep{fashi_dr2}, galaxies with another galaxy inside their flux aperture carry
more H\,I at fixed optical luminosity than isolated ones (\code{fashienv},
\citealt{janskyresearch}). The excess could be blending or real gas. \citet{haynes2018}
describe confusion in ALFALFA as relatively minor (their arXiv version, p.~15, citing
\citealt{jones2016}), and \citet{jones2015} studied it by simulation; we know of no test on the
same galaxies with two beams. FASHI DR2 gives a FAST beam of \hbRealBeamFast~arcmin at
1420~MHz (its Table~1), and \citet{giovanelli2005} an ALFA beam averaging
\hbRealBeamAlfa~arcmin (arXiv preprint, pp.~8 and 22; 3.3 by 3.8~arcmin). We use these nominal
widths; broader effective beams would change only the scale of the expected slope. FASHI DR2 (\hbRealNFashi\ sources) was matched to ALFALFA $\alpha$.100 code-1
detections \citep[\hbRealNAlfalfa\ sources in all;][]{haynes2018} within 1.5~arcmin and
100~km~s$^{-1}$; shifting one catalogue by 10~arcmin gives \hbRealChance\% chance matches.

A neighbour of flux $S_n$ at separation $s_n$, overlapping in velocity, adds
$S_n\exp(-4\ln2\,s_n^2/\theta^2)$ to a beam of width $\theta$. The predicted log ratio $R$ of the
ALFA and FAST fluxes peaks at \hbRealSignal~dex for one equal-flux neighbour at 2~arcmin; its
median is \hbRealRprimMed\ (primary) and \hbRealRnullMed\ (null). The statistic is the least-squares slope $\beta$ of
the observed log flux ratio on $R$, after removing the survey-to-survey flux scale: blending
predicts $\beta\approx1$, extra gas in close pairs $\beta\approx0$. The flux scale (a function of
signal-to-noise, line width, flux and declination) was fitted on the \hbRealNIsolated\ targets
with no neighbour within 6~arcmin and applied to the \hbRealNPrimary\ with a neighbour at
1--6~arcmin in the same velocity window. Here signal-to-noise is the noise-weighted geometric
mean of the two fluxes over the FASHI flux error, not the catalogue SNR. The negative control
uses \hbRealNNull\ targets whose neighbours lie more than 600~km~s$^{-1}$ away and cannot blend;
its slope must be within $2\sigma$ of zero. Errors come from a bootstrap over groups of targets
linked within 6~arcmin.

\section{Results}
\label{sec:results}
On simulated fluxes with the real geometry, a planted $\beta=1$ was detected at $3\sigma$ in
20 of 20 trials. Blends injected into real isolated targets returned $\beta=\hbRealInjOne$
(scatter \hbRealInjOneSd\ over ten injections), within the pre-set tolerance of 0.3; we later
found that the injected flux also raises the calibration covariates, which inflates $\beta$;
with covariates taken before injection the gain is \hbRealConeFrozen. The negative control failed,
$\beta_{\rm null}=\hbRealBetaNull\pm\hbRealBetaNullErr$, so the primary slope
($\hbRealBeta\pm\hbRealBetaErr$) cannot be interpreted.

Diagnostics run afterwards, with their expected outcomes recorded first, did not find the cause.
Tested on isolated targets it was not fitted on, the calibration's bin means are off by up to
\hbRealVoneMaxMean~dex (Figure~\ref{fig:heldout}). Mapped onto the null targets by their
signal-to-noise, this misfit predicts $\beta_{\rm null}=\hbRealMisfitNull\pm\hbRealMisfitNullErr$.
Resampling 2--8-degree strips of right ascension instead of 6-arcmin groups leaves the failure at
\hbRealStripSigMin--\hbRealStripSigMax$\sigma$, and local galaxy density does not account for it.

A second attempt, designed after these diagnostics on the same targets but fixed before it was
run, used a linear spline in signal-to-noise and added local density. It placed a held-out gate
first: fitted on alternate 2-degree strips and tested on the rest, every decile median had to
lie within 0.010~dex of zero with $p>0.01$ over the ten. Both directions failed ($p=\hbRealVtwoPA$
and \hbRealVtwoPB), and the run stopped before $\beta$ was measured. The gate was poorly
chosen. Given the decile errors, a perfect calibration passes the 0.010~dex rule in both
directions with probability \hbRealAmpPass. Medians also respond to skew, while the fit uses
means: in bin means the spline calibration gives $p=\hbRealMeanPA$ and \hbRealMeanPB\
(Figure~\ref{fig:heldout}), and the original form $p<10^{-5}$.

\begin{figure}
\plotone{figures/hiblend_heldout.pdf}
\caption{Mean flux-ratio residual of isolated targets not used in the fit, against
signal-to-noise. The calibration was fitted on alternate 2-degree strips of right ascension and
tested on the others (filled and open symbols for the two directions). Grey: the original
calibration; blue: the spline calibration of the second attempt.}
\label{fig:heldout}
\end{figure}

\section{Conclusion}
\label{sec:conclusion}
The test is inconclusive. Its negative control fails for a reason we have not found: not the
signal-to-noise calibration misfit, not correlated errors on the scales we tried, and not local
density. Misfit along another covariate, or confusion in the broader gridded ALFALFA beam, is
untested. The second attempt's calibration is adequate in the mean, but its gate stopped the
run. A further attempt needs a gate whose false-failure rate is computed before it is fixed,
applied to the statistic the estimator uses. Interferometric H\,I imaging of the pairs would
settle the question directly.

Reproducibility: \code{jansky\_research.hiblend} and \code{scripts/hiblend\_*.py}; every number
here is generated from \code{results/hiblend\_*.json}. The plans, with the controls as frozen,
and the record of each step are \code{plans/97}, \code{plans/98} and
\code{survey/hiblend-findings.md}. Both catalogues are public.

\section*{AI use disclosure}
This paper was developed collaboratively by the author with Anthropic's Claude
\citep{anthropicclaude} (model Claude Opus~5.5, accessed 2026). Claude contributed to the
design, implementation, and testing of the \code{hiblend} module and the drafting of this
manuscript. The author directed the work, independently reviewed and verified all code, results,
and citations, and takes full responsibility for the content. An AI/LLM is not an author.

\begin{acknowledgments}
We used the FASHI DR2 catalogue (FAST, National Astronomical Observatories, Chinese Academy of
Sciences) and the ALFALFA $\alpha$.100 catalogue through VizieR; the analysis builds on the
\code{jansky} project. The author acknowledges the collaborative contribution of Anthropic's
Claude.
\end{acknowledgments}

\software{\code{jansky-research} \citep{janskyresearch}, Claude \citep{anthropicclaude},
astroquery, Astropy, NumPy, SciPy, Matplotlib, jansky \citep{jansky}}

\bibliography{refs}
\end{document}
